% !TeX root = writeup.tex

	\section{Characterization of Fairness Solutions\label{sec:char_fairness_solns}}
	


\subsection{Derivative Computation for $\eqop$\label{sec:char_fairness_solns_eqopp}}
In this section, we prove Lemma~\ref{lem:bijection_tpr_accrate}, which we recall below.
\bijectiontpr*
We will prove Lemma~\ref{lem:bijection_tpr_accrate} in tandem with the following derivative computation which we applied in the proof of Theorem~\ref{thm:eqop_selection}.

\begin{lem}\label{lem:const_to_util_derv}
	The function
	\begin{eqnarray*}
	\Util\supj(t;\boldwj) := \Util\supj\left(\ratef_{\distj}^{-1}\left(\constfj^{-1}(t )\right)\right)
	\end{eqnarray*}
	is concave in $t$ and has left and right derivatives
	\begin{eqnarray*}
	\partial_+ \Util\supj(t;\boldwj) =  \frac{\util(\RCDFj(\constfj^{-1}(t)))}{\boldwj(\RCDFj(\constfj^{-1}(t)))} 
	&\text{and}& \partial_- \Util\supj(t;\boldwj) = \frac{\util(\RCDFplj(\constfj^{-1}(t)))}{\boldwj(\RCDFplj(\constfj^{-1}(t)))}\:.
	\end{eqnarray*}
	\end{lem}

\begin{proof}[Proof of Lemmas~\ref{lem:bijection_tpr_accrate}~and~\ref{lem:const_to_util_derv}] Consider a $\accrate \in [0,1]$. Then, $\distj \circ \ratef_{\distj}^{-1}(\accrate)$ is continuous and left and right differentiable by Lemma~\ref{lem:der_comp}, and its left and right derivatives are indicator vectors $\bolde_{\RCDFj(\accrate)}$ and $\bolde_{\RCDFplj(\accrate)}$, respectively. 
	Consequently, $\accrate \mapsto \langle \boldwj, \distj \circ \ratef_{\distj}^{-1}(\accrate) \rangle$ has left and right derivatives $\boldwj(\RCDF(\accrate))$ and $\boldwj(\RCDFpl(\accrate))$, 
	respectively; both of which are both strictly positive by the assumption $\boldw(x) > 0$. 
	Hence, $\constfj(\accrate) = \langle \boldwj, \distj \circ \ratef_{\distj}^{-1}(\accrate) \rangle$ is strictly increasing  in $\accrate$, and so the map is injective. 
	It is also surjective because $\accrate = 0$ induces the policy $\polj = \mathbf{0}$ and $\accrate = 1$ induces the policy $\polj = \mathbf{1}$ (up to $\distj$-measure zero). Hence, $\constfj(\accrate)$ is an order preserving bijection with left- and right-derivatives, and we can compute the left and right derivatives of its inverse as follows:
	\begin{eqnarray*}
	\partial_+ \constfj^{-1}(t) = \frac{1}{\partial_+ \constfj(\accrate) \big{|}_{\accrate = \constfj^{-1}(t)}} = \frac{1}{\boldwj(\RCDFj(\constfj^{-1}(t)))}~,
	\end{eqnarray*}
	and similarly, $\partial_- \constfj^{-1}(t) = \frac{1}{\boldwj(\RCDFpl(\constfj^{-1}(t)))}$. Then we can compute that
	\begin{eqnarray*}
	\partial_+ \Util\supj(\ratef_{\distj}(\constfj^{-1}(t))) &=& \partial_+\Util(\ratef_{\distj}(\accrate)) \big{|}_{\accrate = \constfj^{-1}(t))} \cdot \partial_+ \constfj(\sup(t))\\
	&=& \frac{\util(\RCDFj(\constfj^{-1}(t)))}{\boldwj(\RCDFj(\constfj^{-1}(t)))}~.
	\end{eqnarray*}
	and similarly $\partial_- \Util\supj(\ratef_{\distj}(\constfj(t))) = \frac{\Util(\RCDFplj(\constfj^{-1}(t)))}{\boldwj(\RCDFplj(\constfj^{-1}(t)))}$. One can verify that for all $t_1 < t_2$, one has that $\partial_+ \Util\supj(\ratef_{\distj}(\constfj^{-1}(t_1))) \ge \partial_- \Util\supj(\ratef_{\distj}(\constfj^{-1}(t_2)))$, and that for all $t$, $\partial_+ \Util\supj(\ratef_{\distj}(\constfj^{-1}(t))) \le \partial_- \Util\supj(\ratef_{\distj}(\constfj^{-1}(t)))$. These facts establish that the mapping $t\mapsto \Util\supj(\ratef_{\distj}(\constfj^{-1}(t)))$ is concave. 
	\end{proof}



